% P15-U1: full source rate proofs and infinite-mean normalization.
\chapter{Strong-Law Rates and Infinite-Mean Normalization}
\label{sec:p15-strong-law-rates-and-normalization}

Changing the denominator of a sum changes both the useful moment
condition and the necessary centering. A finite second moment gives a
logarithmically enlarged square-root rate. Moments below order two give
the Marcinkiewicz--Zygmund rates. With an infinite first absolute moment,
Feller's theorem instead distinguishes zero normalized sums from
unbounded normalized absolute values through an exact tail series.

The source is Durrett's author Version 5 dated 11 January 2019,
Theorems 2.5.11--2.5.13 and their complete proofs
\cite[pp.~87--90]{DurrettPTEAuthor2019}. The source's
Marcinkiewicz--Zygmund statement concerns \(1<p<2\).
The complementary \(0<p<1\) case, the \(p=1\) case and the moment
converse in Exercise 2.5.2 are supplied separately, so that the
parameter ranges and endpoints are explicit. The arguments use the
already proved independent-series and Kronecker theorems.
The power-moment identity used in the truncation calculations is
Lemma 2.2.13 \cite[pp.~63--64]{DurrettPTEAuthor2019}.

\section{Moment tails and the finite-variance rate}
\label{sec:p15-rate-moment-tools}

\begin{hbdefinition}{def:p15-normalized-sums}{Partial sums, centering and positive denominators}
On a probability space \((\Omega,\mathcal F,\mathbb P)\), let
\(X_1,X_2,\ldots\) be real random variables and put
\[
 S_0=0,\qquad S_n=\sum_{j=1}^nX_j,\qquad
 S_n^{(c)}=S_n-nc\quad(c\in\mathbb R).
\]
A zero rate with positive deterministic denominators \(a_n\) means
\(S_n^{(c)}/a_n\to0\) almost surely. All limits use finite real
summands, although expectations and limit superiors may be infinite.
The logarithm is natural. A displayed \(\log n\) denominator starts
at \(n=2\); changing finitely many positive denominators does not
change the asserted limit.
\end{hbdefinition}

\begin{hblemma}{thm:p15-power-tail-tools}{proved=annex:p15.power-tail-tools}{Power moments and elementary summation bounds}
For a nonnegative extended real random variable \(W\) and \(p>0\),
\[
 \mathbb EW^p=\int_0^\infty p y^{p-1}\mathbb P(W>y)\,dy,
 \qquad
 \sum_{n\geq1}\mathbb P(W^p>n)
       \leq\mathbb EW^p
       \leq1+\sum_{n\geq1}\mathbb P(W^p>n).
\]
The identities and inequalities allow infinite values. The value of
the integral's integrand at zero is taken to be zero.
For \(q>1\) and integer \(N\geq1\),
\[
 \sum_{m=N}^\infty m^{-q}
       \leq\frac{q}{q-1}N^{1-q};
\]
for \(N\geq2\), also
\(\sum_{m=N}^\infty m^{-q}\leq (N-1)^{1-q}/(q-1)\).
If \(p>1\), then
\(\sum_{m=1}^n m^{1/p-1}\leq p n^{1/p}\).
Finally, for real \(u,v\) and \(p>0\),
\[
 |u+v|^p\leq2^p(|u|^p+|v|^p).
\]
Thus adding a finite constant preserves finiteness of a positive-order
absolute moment.
\end{hblemma}

\begin{hbtheorem}{thm:p15-finite-variance-log-rate}{proved=annex:p15.finite-variance-log-rate}{The source finite-variance rate}
Let \(X_n\) be i.i.d., with \(\mathbb EX_1=0\) and
\(\mathbb EX_1^2=\sigma^2<\infty\), where \(\sigma\geq0\). For every
\(\varepsilon>0\),
\[
 \frac{S_n}{n^{1/2}(\log n)^{1/2+\varepsilon}}
                       \longrightarrow0\quad\text{almost surely}.
\]
Zero variance is included. For an i.i.d. sequence with finite second
moment and mean \(\mu\), the same rate applies to \(S_n-n\mu\).
The positive \(\varepsilon\) makes the variance series used in this
proof summable; this result does not assert optimality of its logarithmic
factor.
\end{hbtheorem}

\section{The Marcinkiewicz--Zygmund parameter ranges}
\label{sec:p15-mz-ranges}

\begin{hbdefinition}{def:p15-power-cutoff}{The power cutoff}
For a fixed \(p>0\), define
\[
 Y_m=X_m\mathbf1_{\{|X_m|\leq m^{1/p}\}},\qquad
 T_n=\sum_{m=1}^nY_m,\qquad \mu_m=\mathbb EY_m.
\]
These are deletion truncations: a value at either boundary is retained
and a value beyond it is replaced by zero. Each truncated variable is
bounded and has finite first and second moments. This convention differs
from the strict cutoff used below in Feller's theorem.
\end{hbdefinition}

\begin{hbtheorem}{thm:p15-mz-source-rate}{proved=annex:p15.mz-source-rate}{The complete source theorem for \(1<p<2\)}
Let \(X_n\) be i.i.d., let \(1<p<2\), and suppose
\(\mathbb EX_1=0\) and \(\mathbb E|X_1|^p<\infty\). Then
\[
 \frac{S_n}{n^{1/p}}\longrightarrow0\quad\text{almost surely}.
\]
The proof gives the following explicit bounds for the power-cutoff variables:
\[
 \begin{aligned}
 \sum_{m\geq1}\mathbb P(X_m\ne Y_m)&\leq\mathbb E|X_1|^p,\\
 \sum_{m\geq1}\frac{\operatorname{Var}(Y_m)}{m^{2/p}}
       &\leq\sum_{m\geq1}\frac{\mathbb EY_m^2}{m^{2/p}}
          \leq\frac{4}{p(2-p)}\,\mathbb E|X_1|^p.
 \end{aligned}
\]
The full interval decomposition and sum/integral interchange are given
in the annex. The separate mean estimate is
\[
 |\mu_m|\leq m^{1/p-1}
       \mathbb E\!\left[|X_1|^p\mathbf1_{\{|X_1|>m^{1/p}\}}\right],
 \qquad
 n^{-1/p}\sum_{m=1}^n\mu_m\longrightarrow0.
\]
Finite second moments of the untruncated variables are not required.
\end{hbtheorem}

\begin{hbtheorem}{thm:p15-mz-subunit-rate}{proved=annex:p15.mz-subunit-rate}{The uncentered range \(0<p<1\)}
Let \(0<p<1\), and let \(X_n\) have a common real law with
\(\mathbb E|X_1|^p<\infty\). Even without independence,
\[
 \sum_{n\geq1}\frac{|X_n|}{n^{1/p}}<\infty
       \quad\text{almost surely},\qquad
 \frac1{n^{1/p}}\sum_{m=1}^n|X_m|\longrightarrow0
       \quad\text{almost surely}.
\]
Consequently \(S_n/n^{1/p}\to0\) almost surely. A first moment and
a centering constant are unnecessary in this range.
\end{hbtheorem}

\begin{hblemma}{thm:p15-power-rate-necessity}{proved=annex:p15.power-rate-necessity}{The necessary moment and centering}
Let \(X_n\) be i.i.d., \(p>0\) and \(c\in\mathbb R\).
If \(S_n^{(c)}/n^{1/p}\to0\) almost surely, then
\(\mathbb E|X_1|^p<\infty\). If \(p\geq1\), it also follows that
\(c=\mathbb EX_1\).
If the \(p\)-th absolute moment is infinite, the event
\(\{S_n^{(c)}/n^{1/p}\to0\}\) has probability zero.
\end{hblemma}

\begin{hbtheorem}{thm:p15-mz-full-classification}{proved=annex:p15.mz-full-classification}{The complete classification below order two}
Let \(X_n\) be i.i.d., let \(0<p<2\), and fix \(c\in\mathbb R\).
\begin{enumerate}
\item If \(0<p<1\), then \(S_n^{(c)}/n^{1/p}\to0\) almost surely
if and only if \(\mathbb E|X_1|^p<\infty\). Every finite \(c\)
gives the same rate.
\item If \(1\leq p<2\), then \(S_n^{(c)}/n^{1/p}\to0\) almost
surely if and only if \(\mathbb E|X_1|^p<\infty\) and
\(c=\mathbb EX_1\).
\end{enumerate}
At \(p=1\) this is exactly the integrable strong law together with
its necessary absolute moment. The endpoint \(p=2\) is not included;
a bounded centered i.i.d. counterexample is proved below.
\end{hbtheorem}

\section{Feller's infinite-mean normalization}
\label{sec:p15-feller-normalization}

\begin{hbdefinition}{def:p15-feller-normalizers}{The normalized growth condition and strict truncation}
A Feller normalization is a sequence of finite positive real numbers
\(a_n\) for which \(a_n/n\) is nondecreasing. This implies
\(a_n\geq n a_1\) and hence \(a_n\to\infty\).
For i.i.d. real \(X_n\), write
\[
 W=|X_1|,\qquad q_n=\mathbb P(W\geq a_n),\qquad
 Q=\sum_{n\geq1}q_n.
\]
For this theorem only, use the strict deletion variables
\[
 Y_n=X_n\mathbf1_{\{|X_n|<a_n\}},\qquad T_n=\sum_{j=1}^nY_j.
\]
Thus the removed event is \(\{|X_n|\geq a_n\}\). Put \(a_0=0\)
for the intervals \([a_{m-1},a_m)\); it is not a denominator.
\index{Feller normalization}\index{infinite mean!normalization}
\end{hbdefinition}

\begin{hbtheorem}{thm:p15-feller-infinite-mean}{proved=annex:p15.feller-infinite-mean}{The exact infinite-mean tail-series dichotomy}
Let \(X_n\) be i.i.d. real random variables with
\(\mathbb E|X_1|=\infty\), and let \(a_n\) be a Feller
normalization. Then almost surely
\[
 \limsup_{n\to\infty}\frac{|S_n|}{a_n}
   =\begin{cases}
      0,&Q<\infty,\\
      +\infty,&Q=\infty.
     \end{cases}
\]
In the finite-tail case \(S_n/a_n\to0\) almost surely and necessarily
\(a_n/n\to\infty\). No centering is performed, and the signed
expectation \(\mathbb EX_1\) need not exist.
The proof retains the source's scaled-tail argument and its exact
interval-count identity
\[
 \sum_{m\geq1}m\,\mathbb P(a_{m-1}\leq W<a_m)
       =\sum_{n\geq1}\mathbb P(W\geq a_{n-1})=1+Q.
\]
It also supplies the full truncated-variance and truncated-mean arguments
for the finite-tail case.
\end{hbtheorem}

\section{Calculations and sharp boundaries}
\label{sec:p15-rate-calculations}

\begin{hbtheorem}{thm:p15-rate-calculated-examples}{proved=annex:p15.rate-calculated-examples}{Three scales and an infinite-mean normalization}
The following sequences are i.i.d.
\begin{enumerate}
\item For fair \(\{-1,1\}\)-valued signs,
\(S_n/(\sqrt n\log n)\to0\) almost surely.
\item A positive law with mass
\((2^{3/4}-1)2^{-3j/4}\) at \(2^j\), \(j\geq1\), satisfies
\[
 \mathbb E\sqrt{X_1}=1+2^{1/4}+\sqrt2<\infty,\qquad
 \mathbb EX_1=\infty.
\]
Thus \(S_n/n^2\to0\) while \(S_n/n\to\infty\), almost surely.
\item The density
\(f(x)=(7/4)x^{-11/4}\mathbf1_{[1,\infty)}(x)\)
is normalized, has mean \(7/3\), \(3/2\)-moment seven and infinite
second moment. Consequently
\((S_n-(7/3)n)/n^{2/3}\to0\) almost surely.
\item For the positive law \(\mathbb P(X_1=2^j)=2^{-j}\),
the choices \(a_n=n\) and \(a_n=n^2\) give respectively
\[
 \limsup_n S_n/n=\infty,\qquad S_n/n^2\longrightarrow0
                                  \quad\text{almost surely}.
\]
The first has a divergent threshold series and the second a convergent
one.
\end{enumerate}
\end{hbtheorem}

\begin{hbtheorem}{thm:p15-rate-hypothesis-boundaries}{proved=annex:p15.rate-hypothesis-boundaries}{Moment, centering, independence and the order-two endpoint}
\begin{enumerate}
\item For independent fair signs,
\[
 \mathbb P(|S_n|/\sqrt n>1/\sqrt2)\geq1/16\quad(n\geq1).
\]
Hence \(S_n/\sqrt n\) does not tend to zero in probability, and
its almost-sure zero-limit event has probability zero. The proposed
Marcinkiewicz--Zygmund zero rate also fails for every \(p>2\),
despite finiteness of every moment.
\item The constant sequence \(X_n=1\) has all moments, but
\(S_n/n^{1/p}=n^{1-1/p}\) is not a zero rate for any \(p\geq1\).
It also fails the uncentered finite-variance logarithmic conclusion.
\item For any \(p>0\), the law with mass
\((2^p-1)2^{-pj}\) at \(2^j\), \(j\geq1\), is normalized and
has infinite \(p\)-th moment. For every finite \(c\), its independent
sums fail the zero \(n^{1/p}\) rate almost surely. When \(p>1\),
its first moment is finite, so centering by the actual mean does not
repair the missing \(p\)-th moment.
\item The infinite-absolute-mean assumption in Feller's theorem is
substantive: \(X_n=1\) and \(a_n=n\) give \(Q=1<\infty\),
but \(S_n/a_n=1\).
\item Independence in Feller's theorem is substantive: if every
\(X_n\) equals the same positive finite random variable \(X\) with
\(\mathbb P(X=2^j)=2^{-j}\), and \(a_n=n\), then \(Q=\infty\)
but \(|S_n|/a_n=X<\infty\) at every outcome.
\end{enumerate}
\end{hbtheorem}

\begin{hbtheorem}{thm:p15-feller-ratio-counterexample}{proved=annex:p15.feller-ratio-counterexample}{Increasing denominators alone do not suffice}
There exist i.i.d. positive finite \(X_n\) with infinite mean and
positive nondecreasing \(a_n\to\infty\) such that
\[
 \sum_{n\geq1}\mathbb P(X_1\geq a_n)<\infty,\qquad
 \limsup_{n\to\infty}S_n/a_n=\infty
                                    \quad\text{almost surely}.
\]
An explicit deterministic construction takes the dyadic law
\(\mathbb P(X_1=2^j)=2^{-j}\), sets \(N_0=0\), and chooses
\(N_k\) to be the least integer \(n\geq N_{k-1}+2\) such that
\(\mathbb P(S_n/n>k^3)\geq1-2^{-k}\). Then set
\[
 a_n=N_k k^2\qquad(N_{k-1}<n\leq N_k).
\]
The proof establishes existence of every \(N_k\), the convergent
threshold series and the asserted limit superior. Within each block
the ratio \(a_n/n\) decreases, so the example retains increasing
denominators while violating the precise ratio condition.
\end{hbtheorem}

\begin{proofboundary}{def:p15-rates-next-boundary}{Weak laws, distributional convergence and the central limit theorem}
The three selected source rate/normalization proofs are complete here.
The direct triangular-array, tail-truncation and integrable weak laws
are supplied in \cref{thm:p15-triangular-array-weak-law,
thm:p15-iid-weak-tail-criterion,thm:p15-direct-integrable-weak-law}.
Feller's distinct median-based necessity argument is
\cref{thm:p15-feller-generalized-weak-law}.
The St.~Petersburg calculation in Example~2.2.16 is supplied by the
separate normalization arguments in
\cref{thm:p15-petersburg-array-calculation,
thm:p15-petersburg-normalizers-and-prices}.
The numbered Section~3.2 weak-convergence presentation is supplied in
\cref{sec:p15-weak-convergence-of-laws}.
Characteristic functions and the continuity theorem in Section~3.3
are the next prerequisites for the i.i.d. and Lindeberg--Feller
central-limit source proofs in Section~3.4.

The source's later comparison with the law of the iterated logarithm is
\[
 \limsup_n\frac{S_n}{\sqrt{n\log\log n}}=\sigma\sqrt2
\]
under the centered i.i.d. finite-variance hypotheses, with the displayed
denominator used for \(n\geq3\)
\cite[p.~87; Theorem~8.5.2]{DurrettPTEAuthor2019}. That separately
owed theorem is not used or proved here. In particular the failure
at the square-root scale does not declare the present logarithmic
factor optimal.
\end{proofboundary}
